\begin{figure}[t]
\centering
\figplaceholder[2.3in]{\columnwidth}{Fig.~3 --- acting on a moving target (schematic)}{%
\textbf{Layout:} one row, three panels (a)--(c) of equal width, clean vector schematic in the paper's semantic colors; the same robot arm and T-shaped part in all panels; no photographs.
\textbf{(a) Same pose, different motion:} two histories of the part, A (\textcolor{obscol}{solid blue}) and B (\textcolor{obscol}{dashed blue}), with ghost poses at $-100$, $-50$, $0$\,ms that end in the same pose at capture $\tau_t$. Label: ``current pose alone is ambiguous''.
\textbf{(b) Pose at execution:} a timeline strip under the panel: capture $\tau_t$, then ticks every $\Delta_a=50$\,ms at the times $\tapp_{t,j}$ of the chunk's commands; the ticks before the chunk arrives at $\tau_t+\ell_t$ are grayed out as expired, and the first executed command is marked $\tapp$. Above it, the part at $\tapp$ for A and B with different positions and rotations (\textcolor{predcol}{orange} triads). Label: ``the command must fit the pose at $\tapp$''.
\textbf{(c) Roles:} one scene with three kinds of bodies. \emph{Target}: the part with the end-effector frame $E$, the target frame $O$, the contact point, and tolerance $\varepsilon$ (mm). \emph{Context, obstacle}: a receptacle wall or occluder with a clearance arrow $m\gg\varepsilon$. \emph{Context, driver}: the conveyor carrying the part, with a motion arrow, and optionally a striker approaching it; label ``moves the target''. \emph{Background}: table edge and distant objects desaturated in \textcolor{gtcol}{gray}; label ``no effect''. A three-row legend: target $\to$ precise SE(3); context $\to$ coarse 3D flow; background $\to$ nothing.
\textbf{Rules:} conceptual illustration with hypothetical positions; \textcolor{obscol}{blue} past, \textcolor{predcol}{orange} future, \textcolor{actcol}{purple} command, \textcolor{gtcol}{gray} background.}
\caption{\textbf{Acting on a moving target} (schematic). (a) Two histories reach
the same pose at capture time $\tau_t$ but imply different motions. (b) At the
execution time $\tapp$, they lead to different target poses. (c) A contact
requires the end-effector frame $E$ to match the target $O$ within $\varepsilon$;
an obstacle in the context matters only through the gap $m\gg\varepsilon$ it leaves,
a driver such as the conveyor matters through the motion it imparts to the target,
and the background does not affect the task.}
\label{fig:principles}
\end{figure}